\chapter{ADR-028 --- Bố cục source theo miền (chốt và áp 2026-10-12)}
\label{chap:adr28}
\sloppy

\emph{\textbf{Trạng thái: đã lên UAT 12/10/2026} (PR \#7 vào \texttt{deploy/staging}, commit \texttt{f3107562}).
Các bước đã làm trên máy chủ ở \S\ref{sec:adr28-deploy}. ADR-027 (chapter 74) trả lời
``module phụ thuộc nhau thế nào''; ADR này trả lời ``module nằm ở thư mục nào''.}

\section{Bối cảnh}

Trước 12/10 repo có \texttt{odoo19/} và một thư mục \texttt{custom/} phẳng chứa 66 module: 44 module Wujia trộn với
22 module bên ngoài (\texttt{om\_*}, \texttt{eh\_*}, OCA, MuK, MCP, dashboard). \texttt{themes/} giữ thêm 15 module,
trong đó 8 module \texttt{queue\_job} của OCA nằm lồng trong \texttt{themes/oca\_queue}. Nhìn cây thư mục không biết
phân hệ nào ở đâu, không tách được code của mình với code người khác, và không có chỗ cho phần riêng của từng nơi
triển khai.

Chủ dự án chốt hướng đi theo cách Beesuite tổ chức: Odoo để riêng, module chia thư mục theo phân hệ, thêm một thư mục
\texttt{addons/} cho phần riêng. Wujia là \textbf{core}; mỗi quốc gia hoặc bên nhượng quyền triển khai bằng một bản
clone cả core, phần riêng bỏ vào \texttt{addons/}.

\section{Quyết định}

\begin{lstlisting}[style=shell]
Odoo19C/          Odoo 19 Community, ngoai repo (hoac odoo19/ trong repo, gitignore)
WujiaTea/
  config/  docs/  scripts/
  base/           wujia_core, wujia_portal_layout, wujia_mobile_core, wujia_i18n, wujia_audit
  franchise/      wujia_franchise, _contract, _inspection, _operations, wujia_info_request
  sales/          wujia_sale, wujia_order_window, wujia_return
  logistics/      wujia_fleet, wujia_delivery
  account/        wujia_account
  academy/        wujia_knowledge, wujia_exam
  service/        wujia_support, wujia_notification
  portal/         wujia_portal_base + 12 wujia_portal_<chuc nang>
  mobile/         7 wujia_mobile_<x>
  bi/             wujia_metabase_connector, wj_ks_dashboard_ninja, wj_ks_dn_advance
  tools/          mcp_server, muk_mcp, muk_web_utils, login_as_any_user, wujia_fields_value
  themes/         7 muk_web_*, legion_enterprise_theme
  third_party/    oca/ (4)  oca_queue/ (8)  om/ (8)  misc/ (5)
  addons/         rong o core, module rieng cua tung ban trien khai
\end{lstlisting}

81 module, 17 mục \texttt{addons\_path}. Tên thư mục viết thường, vì máy chủ Linux phân biệt hoa/thường còn máy Mac
thì không. \texttt{UI/} (file thiết kế) dời vào \texttt{docs/design/}.

\subsection{Luật đặt module}

\begin{enumerate}[nosep]
\item Thư mục là \textbf{miền nghiệp vụ}; tầng module vẫn theo ADR-027. \texttt{base/} chứa nền và hai khung kênh,
  \texttt{portal/} và \texttt{mobile/} chứa module ghép, các thư mục còn lại chứa module nghiệp vụ.
\item \texttt{wujia\_portal\_<x>} luôn ở \texttt{portal/}, \texttt{wujia\_mobile\_<x>} luôn ở \texttt{mobile/}, kể cả
  khi nghiệp vụ của nó nằm ở thư mục khác.
\item \texttt{third\_party/} và \texttt{themes/} là code bên ngoài: muốn đổi hành vi thì kế thừa ở module Wujia.
  Ngoại lệ đang có: \texttt{dynamic\_dashboard\_ai\_nexgen} và \texttt{dynamic\_approval\_workflow} đã bị vá tại
  chỗ (vá SQL injection); sửa tại chỗ thì ghi rõ trong commit.
\item \texttt{addons/} ở core luôn rỗng. Bản clone chỉ thêm module vào \texttt{addons/}, không sửa thư mục core;
  thiếu điểm mở thì sửa ở core rồi kéo về.
\item Thêm thư mục mới: khai trong \texttt{ADDON\_DIRS} của \texttt{scripts/wj\_paths.py}, rồi sinh lại
  \texttt{addons\_path} ở mọi máy.
\end{enumerate}

Luật 1 và 2 được máy kiểm: \texttt{scripts/qa/check\_layers.py} thêm luật \textbf{R8} (module đặt đúng thư mục của
tầng nó). Kết quả sau khi dời: 0 vi phạm tầng, 0 vi phạm R8.

\subsection{Vì sao gom theo miền, không mỗi phân hệ một thư mục}

Nhiều phân hệ hiện chỉ có một module (công nợ, hỗ trợ, thông báo). Mỗi phân hệ một thư mục sẽ ra gần 20 thư mục một
module. Gom theo miền cho 10 thư mục Wujia, mỗi thư mục 1--13 module, và vẫn đọc ra phân hệ. \texttt{portal/} và
\texttt{mobile/} đứng riêng vì chúng là \emph{kênh}: một màn portal ghép nhiều nghiệp vụ, không thuộc riêng miền nào.

\subsection{Mô hình clone}

Bản triển khai mới = clone repo core, thêm module vào \texttt{addons/}, sinh \texttt{addons\_path} bằng
\texttt{wj\_paths.py}. Module trong \texttt{addons/} được kế thừa hoặc depend mọi module core; core không bao giờ
depend ngược lại. Trên máy chủ đã có một bản clone kiểu này: \texttt{huipeng.hongtrangogia.com} (repo
\texttt{huifeng.erp.staging}), còn ở bố cục cũ \texttt{custom/} + \texttt{themes/}.

\section{Một nguồn đường dẫn}

\texttt{scripts/wj\_paths.py} là chỗ duy nhất biết bố cục thư mục.

\begin{lstlisting}[style=shell]
python3 scripts/wj_paths.py --addons-path            # dong addons_path cho odoo.conf
python3 scripts/wj_paths.py --where wujia_exam       # thu muc cua mot module
python3 scripts/wj_paths.py --list                   # folder -> module
python3 scripts/wj_paths.py --changed-modules A B    # module doi noi dung giua 2 commit
\end{lstlisting}

Mười script trong \texttt{scripts/qa/}, \texttt{i18n\_tool.py}, \texttt{sync\_translations.py} và bốn script shell
(\texttt{start}, \texttt{upgrade}, \texttt{init-db}, \texttt{reseed\_full}) đã chuyển sang hỏi \texttt{wj\_paths}
thay vì ghép chuỗi \texttt{custom/<module>}. Odoo core được dò theo thứ tự: biến \texttt{WUJIA\_ODOO}, rồi
\texttt{../Odoo19C}, rồi \texttt{./odoo19}; dời Odoo ra ngoài repo lúc nào cũng được, không phải sửa script.

\texttt{--changed-modules} bỏ qua file chỉ đổi chỗ. Thiếu điều này thì lần deploy đầu sau khi dời sẽ thấy cả 81
module ``có thay đổi'' và nâng cấp tất cả.

\section{Vì sao dời mà không phải migrate}

Tên kỹ thuật module không đổi. Vì vậy \texttt{ir\_module\_module}, \texttt{ir\_model\_data}, đường dẫn view trong
DB, URL \texttt{/<module>/static/\ldots}, thư mục \texttt{migrations/} và file dịch đều không đổi. Odoo tìm module
bằng cách quét các thư mục trong \texttt{addons\_path}, nên việc duy nhất phải làm trên máy chạy là đổi dòng
\texttt{addons\_path} rồi khởi động lại.

\section{Phát hiện trong lúc làm}

\begin{enumerate}[nosep]
\item \textbf{30 file test đọc source bằng thư mục cha.} Các test ``quét source'' (kiểm CSS, QWeb của mọi màn portal)
  tự tính thư mục \texttt{custom/} rồi đọc file của module bên cạnh. Sau khi dời, 72 test đỏ, và tệ hơn là vài test
  khác \emph{xanh giả}: vòng quét không còn thấy module nào nên không kiểm gì. Đã thêm
  \texttt{wujia\_portal\_layout/tests/src\_paths.py} (hỏi Odoo module nằm ở đâu) và chuyển cả 30 file sang dùng nó.
  Chỉ đổi cách tìm đường dẫn, không đổi điều test kiểm.
\item \textbf{Khung portal import ngược lên \texttt{wujia\_portal\_base}.} \texttt{wujia\_portal\_layout} lấy
  \texttt{rate\_limit} và \texttt{build\_pager} từ \texttt{portal\_base}, trái luật 2 ADR-027 (khung không phụ thuộc
  module ghép). Đã dời hai helper này xuống khung; \texttt{portal\_base} import lại đúng tên cũ nên các controller
  đang gọi không đổi dòng nào.
\item \textbf{\texttt{wujia\_mobile\_core} depend \texttt{base, web}}, trong khi ADR-027 ghi \texttt{web} +
  \texttt{muk\_web\_theme}. Chỉ ghi nhận; đây là code anh Thái.
\item \textbf{\texttt{deploy.sh} trên máy chủ dò module đổi bằng \texttt{\^{}(custom|themes)/}.} Sau khi dời, phép dò
  này luôn trả rỗng. Phải sửa một dòng trong file đó (\S\ref{sec:adr28-deploy}).
\end{enumerate}

\section{Kiểm chứng}

\begin{center}
\small
\begin{tabularx}{\textwidth}{>{\raggedright\arraybackslash}p{5.2cm}>{\raggedright\arraybackslash}X}
\toprule
\textbf{Phép kiểm} & \textbf{Kết quả} \\
\midrule
Nội dung 81 module trước và sau khi dời (tree SHA của git) & giống nhau cả 81 \\
Cùng một lệnh test, hai bản sao DB: bố cục cũ và bố cục mới &
  cả hai: 1359 test, 0 đỏ (chạy lại sau khi sửa test và vá khung: vẫn 1359, 0 đỏ) \\
Số test theo từng module (21 module) & giống nhau từng module \\
File source mà các test quét nhìn thấy & 2092 file ở cả hai bố cục (bố cục mới thêm đúng 1 file helper) \\
\texttt{check\_layers.py} & 0 vi phạm tầng, 0 vi phạm R8 \\
Kết quả các script QA trước và sau & giống nhau; \texttt{controller\_inventory -{}-all} nay thấy thêm 8 module
  \texttt{queue\_job} \\
Khởi động server trên bố cục mới & \texttt{/web/login}, \texttt{/portal/login} trả 200; file tĩnh từ \texttt{base/}, \texttt{themes/}, \texttt{third\_party/oca\_queue}, \texttt{third\_party/om} trả 200; 92 module \texttt{installed} ở cả hai bản \\
\bottomrule
\end{tabularx}
\end{center}

Bộ test: \texttt{-u} 21 module (khung, \texttt{portal\_base}, 11 module portal, 7 module nghiệp vụ khối A,
\texttt{wujia\_i18n}) kèm \texttt{-{}-test-enable}, chạy trên bản sao của DB \texttt{wujia\_vr}. Chưa đo được trên
UAT vì chưa có tài khoản UAT mới.

\section{Máy chủ: hiện trạng và các bước đổi}
\label{sec:adr28-deploy}

Khảo sát chỉ đọc qua aaPanel ngày 12/10/2026. Không sửa, không khởi động lại gì.

\begin{center}
\small
\begin{tabularx}{\textwidth}{>{\raggedright\arraybackslash}p{3.3cm}>{\raggedright\arraybackslash}X}
\toprule
\textbf{Mục} & \textbf{Hiện trạng} \\
\midrule
Source & \texttt{/www/wwwroot/dev.wujiatea.com.vn}, nhánh \texttt{deploy/staging}, commit \texttt{ded91f1b} \\
Tiến trình & systemd \texttt{odoo-dev}, user \texttt{www}, cổng 8019 (websocket 8020), nginx đứng trước \\
File cấu hình & \texttt{config/odoo.conf} trong thư mục source (không nằm trong git) \\
\texttt{addons\_path} hiện tại & \texttt{odoo19/addons}, \texttt{themes}, \texttt{themes/oca\_queue}, \texttt{custom},
  và một thư mục ngoài repo: \texttt{/www/wwwroot/wujia-uat.tipscode.io/odoo\_modules} \\
Deploy & bấm \emph{Deploy latest} trong Git Manager của site (Automatic deployment: Off); aaPanel kéo code, git hook \texttt{post-merge} gọi \texttt{deploy.sh -{}-no-pull}: dò module đổi, chạy
  \texttt{-u}, khởi động lại, gọi \texttt{/web/login}, báo Discord \\
Hai instance còn lại & \texttt{odoo-uat} (DB \texttt{wujia\_mn}, cổng 8029, repo khác, thư mục
  \texttt{odoo\_modules}) và \texttt{odoo-huipeng} (DB \texttt{huipeng\_db}, cổng 8059, repo
  \texttt{huifeng.erp.staging}). Ba tiến trình riêng, ba thư mục source riêng \\
\bottomrule
\end{tabularx}
\end{center}

Hai điều rút ra. Một: khởi động lại \texttt{odoo-dev} không đụng hai instance kia. Hai: \textbf{merge PR không tự deploy}; phải bấm \emph{Deploy latest}. Nếu lúc bấm \texttt{addons\_path} còn trỏ
\texttt{custom/} thì Odoo khởi động lại mà không thấy module Wujia nào, nên hai chỗ dưới đây phải sửa trước.

\subsection{Việc đã làm trên máy chủ (12/10/2026), theo thứ tự}

\begin{enumerate}[nosep]
\item Sửa dòng \texttt{addons\_path} trong
  \texttt{/www/wwwroot/dev.wujiatea.com.vn/config/odoo.conf}: thêm 16 thư mục mới vào, \emph{giữ nguyên} các mục
  cũ. Odoo bỏ qua thư mục chưa tồn tại (chỉ ghi cảnh báo), nên dòng này đúng cả trước lẫn sau khi merge, và không cần
  khởi động lại ở bước này. Giá trị đầy đủ nằm trong mô tả PR.
\item Sửa dòng dò module trong \texttt{deploy.sh} (dòng 222) thành:
  \texttt{CHANGED\_MODULES=\$("\$PYTHON\_BIN" scripts/wj\_paths.py -{}-changed-modules "\$OLD\_HEAD" "\$NEW\_HEAD")}
\item Merge PR \#7, bấm \emph{Deploy latest}. Máy chủ từ \texttt{ded91f1b} lên \texttt{f3107562}, \texttt{odoo-dev} khởi động lại.
  \texttt{deploy.sh} \emph{không} chạy \texttt{-u}: ở đường này \texttt{OLD\_HEAD = NEW\_HEAD} và HEAD là merge commit, nên
  dò ra rỗng. PR này chỉ đổi Python, \texttt{.po} và hai số version nên không cần \texttt{-u}; PR có đổi view, data,
  security hay migration thì phải \texttt{-u} tay cho tới khi \texttt{deploy.sh} lấy \texttt{ORIG\_HEAD} làm commit cũ.
\item Đã kiểm từ ngoài: \texttt{/web/login} và \texttt{/portal/login} trả 200, file tĩnh của 13 thư mục trả 200, log sau
  khởi động lại không có ERROR, module nạp từ đường dẫn mới.
\item Còn lại: \texttt{custom/} và \texttt{themes/oca\_queue} trên máy chủ chỉ còn vỏ \texttt{\_\_pycache\_\_}; xoá rồi bỏ hai
  mục đó khỏi \texttt{addons\_path}.
\end{enumerate}

Lùi lại: trả \texttt{odoo.conf} và \texttt{deploy.sh} về bản sao lưu (đuôi \texttt{.bak-adr028}, nằm cạnh file), \texttt{git reset} về commit cũ, khởi động lại.
Không có gì trong DB phải lùi.

Hai nhánh đang mở khác (\texttt{jules}, việc dở của anh Thái) phải merge \texttt{deploy/staging} về sau khi PR vào.
Git nhận ra file đổi chỗ nên sửa nội dung vẫn merge được; riêng file \emph{mới thêm} vào \texttt{custom/} phải dời
tay sang thư mục mới.

\section{Ngoài phạm vi, để lại}

\begin{itemize}[nosep]
\item Dời thật \texttt{odoo19/} ra \texttt{../Odoo19C} (script đã hỗ trợ cả hai).
\item Chuyển bản clone Huipeng sang bố cục mới.
\item Gỡ module bên ngoài không dùng: cần đo trạng thái cài trên UAT trước.
\item \texttt{server\_wide\_modules} trên máy chủ là \texttt{base,web}, máy dev là \texttt{base,web,queue\_job}:
  chưa đối chiếu.
\item \texttt{scripts/qa/css\_move/}: script một lần, còn đường dẫn máy cũ, không sửa.
\item Sửa depend của \texttt{wujia\_mobile\_core}; viết lại CI cho máy aaPanel; chuyển repo sang private.
\end{itemize}

\section{Kết quả nghiệp vụ của phiên}

Người dùng không thấy thay đổi nào. Thứ thay đổi là cách tổ chức mã nguồn: mở repo là thấy phân hệ, code của Wujia
tách khỏi code mua và code cộng đồng, và có sẵn chỗ để mỗi nước hoặc mỗi bên nhượng quyền thêm phần riêng mà không
sửa vào lõi. Việc dời đã được chứng minh là không đổi nội dung và không đổi kết quả test.
